\chapter{空间中的向量、直线与平面}\label{ch:g12:space}

一旦向量与坐标可用，三维几何便可化为计算问题：
直线可用参数方程表示，平面可用笛卡儿方程表示，
数量积则用来度量角与距离。
本章建立这一工具，
并用它求解经典的交点与距离问题。

\section{空间中的向量}

空间中的向量与平面中的向量遵循相同的运算规则：可以相加，可以数乘，
且 $\vect{AB} = \vect{CD}$ 当且仅当 $ABDC$ 为平行四边形。
在坐标系 $(O; \vec\imath, \vec\jmath, \vec k)$ 中，
每个向量都有坐标 $(x, y, z)$。

\begin{definition}[共线与共面]\label{def:g12:space:collinear}
若两个向量中一个是另一个的倍数，则称二者\emph{共线}。
三个向量 $\vec u, \vec v, \vec w$ 若其中一个可以写成另外两个的组合，则称它们\emph{共面}\index{共面}，
例如
$\vec w = a\vec u + b\vec v$。
\end{definition}

\begin{proposition}\label{prop:g12:space:basis}
若 $\vec u, \vec v, \vec w$ 不共面，则空间中的每个向量都可唯一地分解为
$x\vec u + y\vec v + z\vec w$。
\end{proposition}

\begin{proof}[证明思路]
过待分解向量的终点作与 $\vec w$ 平行的直线，
该直线与 $\vec u, \vec v$ 所在平面交于唯一一点，由此把向量分解为
平面内的分量（唯一地为 $x\vec u + y\vec
v$，由平面几何可知）与沿 $\vec w$ 方向的分量。
唯一性：两种分解之差会把其中一个向量用另外两个表示，
这与三者不共面相矛盾。
\end{proof}

\section{直线与平面}

\begin{definition}[直线的参数方程]\label{def:g12:space:line}
过点 $A(x_A, y_A, z_A)$、方向向量为
$\vec u (a, b, c) \neq \vec 0$ 的直线是如下点的集合：
\[
M(x_A + ta,\; y_A + tb,\; z_A + tc), \qquad t \in \R .
\]
\end{definition}

\begin{definition}[平面]\label{def:g12:space:plane}
过点 $A$、以两个不共线的向量 $\vec u, \vec v$ 为方向的平面
是满足下式的点 $M$ 构成的集合：
$\vect{AM} = s\vec u + t\vec v$，$(s, t) \in \R^2$。
\end{definition}

\begin{omfigure}
\begin{tikzpicture}[scale=1.05]
  \draw (-0.5,-0.25) -- (3.4,1.7);
  \draw[-Stealth, omDef, thick] (1,0.5) -- (1.9,0.95)
    node[midway, above left, font=\small] {$\vec u$};
  \fill (0.1,0.05) circle (1.2pt) node[below right, font=\small] {$t=-1$};
  \fill (1,0.5) circle (1.2pt) node[above left, font=\small] {$A$};
  \fill (1.9,0.95) circle (1.2pt) node[below right, font=\small] {$t=1$};
  \fill (2.8,1.4) circle (1.2pt) node[below right, font=\small] {$t=2$};
\end{tikzpicture}
\qquad\quad
\begin{tikzpicture}[scale=1.05]
  \draw (0,0) -- (4,0.5) -- (5.2,1.9) -- (1.2,1.4) -- cycle;
  \node[font=\small] at (4.75,1.6) {$\mathcal P$};
  \coordinate (A) at (1.3,0.55);
  \coordinate (M) at (3.77,1.38);
  \draw[-Stealth, omDef, thick] (A) -- ++(1.31,0.16)
    node[midway, below, font=\small] {$\vec u$};
  \draw[-Stealth, omDef, thick] (A) -- ++(0.42,0.49)
    node[midway, left, font=\small] {$\vec v$};
  \draw[black, dashed, thin] (3.27,0.79) -- (M) -- (1.80,1.14);
  \fill (A) circle (1.2pt) node[below left, font=\small] {$A$};
  \fill (M) circle (1.2pt) node[above, font=\small] {$M$};
\end{tikzpicture}

{\small 左图：过 $A$、方向为 $\vec u$ 的直线，并以参数 $t$ 标出刻度。
右图：过 $A$、以 $\vec u$ 与 $\vec v$ 为方向的平面，
平面上每点 $M$ 均可表示为
$\vect{AM} = s\vec u + t\vec v$。}
\end{omfigure}

\begin{proposition}[相对位置]\label{prop:g12:space:positions}
两个不同的平面要么平行，要么相交于一条直线。
直线与平面要么平行（可能包含于平面内），要么相交于唯一一点。
空间中的两条直线可能相交、严格平行、重合，
或\emph{异面}（不共面）。
\end{proposition}

\begin{proof}[思路简述]
这些都是关联性陈述，每种情形的讨论在坐标下都归结为
解线性方程组并计算其解的个数 --- 见
\cref{met:g12:space:intersections}。
\end{proof}

\section{数量积}

\begin{definition}[数量积]\label{def:g12:space:dot}
向量 $\vec u$ 与 $\vec v$ 的\emph{数量积}\index{数量积}为
\[
\vec u \cdot \vec v = \tfrac12\left(\norm{\vec u + \vec v}^2
- \norm{\vec u}^2 - \norm{\vec v}^2\right),
\]
其中 $\norm{\vec u}$ 为 $\vec u$ 的长度。若两个向量均非零，
则 $\vec u \cdot \vec v = \norm{\vec u}\,\norm{\vec v}\cos\theta$，
其中 $\theta$ 为二者的夹角；在\emph{标准正交}
坐标系中，
\[
\vec u \cdot \vec v = xx' + yy' + zz' .
\]
若两个向量的数量积为 $0$，则称它们\emph{正交}。
\end{definition}

\begin{proposition}[性质]\label{prop:g12:space:dotprops}
数量积是对称的（$\vec u\cdot\vec v = \vec v\cdot\vec u$），
双线性的（$(a\vec u + b\vec v)\cdot \vec w = a\,\vec u\cdot\vec w +
b\,\vec v\cdot \vec w$），并且
$\vec u \cdot \vec u = \norm{\vec u}^2 \geq 0$.
\end{proposition}

\begin{proof}
在标准正交系中，三条性质均可直接由公式
$xx' + yy' + zz'$ 读出，而该公式本身由定义式
及长度的勾股计算得到：
$\norm{\vec u}^2 = x^2 + y^2 + z^2$.
\end{proof}

\begin{definition}[法向量]\label{def:g12:space:normal}
平面 $\mathcal P$ 的\emph{法向量}\index{法向量}是指一个
非零向量，它与 $\mathcal P$ 中每个向量都正交（只需它与
$\mathcal P$ 的两个不共线方向正交，
即可）。
\end{definition}

\begin{omfigure}
\begin{tikzpicture}[scale=1.1]
  \draw (0,0) -- (4,0.5) -- (5.2,1.9) -- (1.2,1.4) -- cycle;
  \node[font=\small] at (4.75,1.6) {$\mathcal P$};
  \coordinate (H) at (2.6,0.95);
  \coordinate (M0) at (2.86,3.03);
  \draw[-Stealth, omThm, thick] (H) -- (2.73,1.99)
    node[midway, right=1pt, font=\small] {$\vec n$};
  \draw[black, dashed, thin] (2.73,1.99) -- (M0);
  \draw[thin] (2.79,0.974) -- (2.815,1.174) -- (2.625,1.15);
  \draw[black, dashed, thin] (H) -- (1.1,0.55);
  \fill (H) circle (1.2pt) node[below right, font=\small] {$H$};
  \fill (M0) circle (1.2pt) node[above, font=\small] {$M_0$};
  \fill (1.1,0.55) circle (1.2pt) node[below left, font=\small] {$A$};
\end{tikzpicture}

{\small 法向量 $\vec n$ 与平面 $\mathcal P$ 的每个方向都
正交。$\mathcal P$ 上的点 $H$ 为由 $M_0$ 向该平面所作垂线的垂足，
它给出 $M_0$ 到平面的距离。}
\end{omfigure}

\begin{theorem}[平面的笛卡尔方程]\label{thm:g12:space:planeeq}
在标准正交坐标系中，过
点 $A$、法向量为 $\vec n(a, b, c)$ 的平面方程为
\[
a x + b y + c z + d = 0,
\]
其中 $d = -(ax_A + by_A + cz_A)$。反之，任一形如上式的方程，
只要 $(a,b,c) \neq (0,0,0)$，就确定一个法向量为 $(a, b, c)$ 的平面。
\end{theorem}

\begin{proof}
$M \in \mathcal P \iff \vect{AM} \perp \vec n \iff
a(x - x_A) + b(y - y_A) + c(z - z_A) = 0$，展开即得该方程。
反之，任取该方程的一个解点 $A$；
同样的计算反推可知，解集为
$\{M : \vect{AM}\cdot\vec n = 0\}$，即过 $A$ 且
以 $\vec n$ 为法向量的平面。
\end{proof}

\begin{method}[交集的实际求法]\label{met:g12:space:intersections}
\begin{itemize}
  \item \emph{直线 $\cap$ 平面}：把直线的参数式坐标
        代入平面方程，解出 $t$（一解：一个
        交点；$0 = $ 非零：平行；$0 = 0$：直线在平面内）。
  \item \emph{平面 $\cap$ 平面}：联立两平面的方程组，
        以一个自由坐标为参数；解为一条直线（若
        两平面的法向量共线，则两平面平行）。
  \item \emph{正交性判别}：直线 $\perp$ 平面当且仅当其方向
        与法向量共线；两平面垂直当且仅当其法向量
        正交。
\end{itemize}
\end{method}

\begin{proposition}[点到平面的距离]\label{prop:g12:space:distance}
\index{距离！点到平面}
在标准正交系中，点 $M_0(x_0, y_0, z_0)$ 到
平面 $\mathcal P : ax + by + cz + d = 0$ 的距离为
\[
\operatorname{dist}(M_0, \mathcal P)
= \frac{\abs{ax_0 + by_0 + cz_0 + d}}{\sqrt{a^2 + b^2 + c^2}}.
\]
\end{proposition}

\begin{proof}
设 $H$ 为 $M_0$ 在 $\mathcal P$ 上的正交投影，即垂线的垂足，
故 $\vect{HM_0}$ 与单位法向量
$\frac{\vec n}{\norm{\vec n}}$ 共线，且距离为
$\abs{\vect{HM_0}\cdot \frac{\vec n}{\norm{\vec n}}}$。对任意
$H \in \mathcal P$，
有 $\vect{HM_0}\cdot\vec n = ax_0 + by_0 + cz_0 - (ax_H + by_H + cz_H)
= ax_0 + by_0 + cz_0 + d$（利用 $H$ 所满足的方程），由此
再除以 $\norm{\vec n} = \sqrt{a^2+b^2+c^2}$ 即得公式。
\end{proof}

\begin{example}\label{ex:g12:space:distance}
原点到平面 $x + 2y + 2z - 6 = 0$ 的距离：
$\frac{\abs{-6}}{\sqrt{1 + 4 + 4}} = \frac63 = 2$.
\end{example}

\section{练习}

所有练习中的坐标系均为标准正交坐标系。

\begin{exercise}[$\star$]\label{exo:g12:space:1}
已知 $A(1, 0, 2)$、$B(3, 1, 1)$、$C(2, -1, 3)$，判断点 $A$、$B$、
$C$ 是否共线，并求直线 $(AB)$ 的参数方程。
\end{exercise}

\begin{exercise}[$\star$]\label{exo:g12:space:2}
求过 $A(1, 1, 0)$ 且法向量为
$\vec n(2, -1, 3)$ 的平面的一般方程，并判断点 $B(0, 2, 1)$ 是否在该平面上。
\end{exercise}

\begin{exercise}[$\star$]\label{exo:g12:space:3}
求该三角形在顶点 $A$ 处的角（精确到
度），该三角形的顶点为 $A(0,0,0)$、$B(1, 1, 0)$ 和 $C(1, 0, 1)$。
\end{exercise}

\begin{exercise}[$\star\star$]\label{exo:g12:space:4}
设直线 $\mathcal D$ 过点 $A(1, 2, 0)$，方向向量为
$\vec u(1, -1, 2)$，又设平面 $\mathcal P : 2x + y - z + 1 = 0$。
求 $\mathcal D \cap \mathcal P$。
\end{exercise}

\begin{exercise}[$\star\star$]\label{exo:g12:space:5}
设 $\mathcal P : x - y + 2z = 1$，$\mathcal Q : 2x + y + z = 4$。
证明 $\mathcal P$ 与 $\mathcal Q$ 相交于一条直线，
并求该直线的参数方程。
\end{exercise}

\begin{exercise}[$\star\star$]\label{exo:g12:space:6}
证明下列两条直线
\[
\mathcal D_1 : (x, y, z) = (1 + t,\; 2t,\; 3 - t), \qquad
\mathcal D_2 : (x, y, z) = (2 + s,\; 1 + s,\; 1 + 2s)
\]
是异面直线，即两者不共面。
\end{exercise}

\begin{exercise}[$\star\star$]\label{exo:g12:space:7}
设 $\mathcal P : 2x - y + 2z - 3 = 0$。
\begin{enumerate}
  \item 求点 $M_0(3, 1, 1)$ 到 $\mathcal P$ 的距离。
  \item 求 $H$ 的坐标：它是
        $M_0$ 在 $\mathcal P$ 上的正交投影，并检验距离 $M_0H$。
\end{enumerate}
\end{exercise}

\begin{exercise}[$\star\star\star$]\label{exo:g12:space:8}
立方体 $ABCDEFGH$ 的棱长为 $1$，建立坐标系，使得
$A(0,0,0)$、$B(1,0,0)$、$D(0,1,0)$、$E(0,0,1)$（其中
$C = B + D$、$F = B + E$、$G = B + D + E$、$H = D + E$ 均为
自 $A$ 出发的向量）。
\begin{enumerate}
  \item 证明对角线 $(AG)$ 与平面 $(BDE)$ 正交。
  \item 求 $A$ 到平面 $(BDE)$ 的距离，以及
        直线 $(AG)$ 与该平面的交点。
\end{enumerate}
\end{exercise}

\begin{exercise}[$\star\star\star$]\label{exo:g12:space:9}
\emph{（四面体的体积。）} 设 $A(1,1,1)$、$B(2,1,0)$、$C(0,2,1)$
和 $D(2,2,2)$。
\begin{enumerate}
  \item 验证向量 $\vect{AB}$、$\vect{AC}$、$\vect{AD}$ 不共面
        （这四点确实构成一个四面体）。
  \item 求向量 $\vec n$，使其与 $\vect{BC}$ 和
        $\vect{BD}$ 均正交，并由此推出平面 $(BCD)$ 的一般方程。
  \item 求 $A$ 到平面 $(BCD)$ 的距离，以及
        三角形 $BCD$ 的面积，利用公式
        $\text{面积} = \frac12\sqrt{\norm{\vec u}^2\norm{\vec v}^2 -
        (\vec u \cdot \vec v)^2}$ 适用于由
        $\vec u, \vec v$ 张成的三角形。
  \item 由此推出四面体的体积
        （$V = \frac13 \times \text{底面积} \times \text{高}$）。
\end{enumerate}
\end{exercise}

\section{问题：切割立方体}

\begin{problem}[{周末问题 --- 一个平面将立方体截出一个
正六边形，一个正四面体藏于角落之中，而对角线的夹角
决定了每一个菱形}]
\label{pb:g12:space:1}
切开立方体，人们原以为会得到正方形和长方形 --- 然而一种
著名的切法会得到一个\emph{正六边形}，另有四个巧妙的
角切割会留下一个\emph{正四面体}。这两个令人惊讶的结果，
以及碳原子所钟情的 $109.5^\circ$ 角，
都可由本章的工具解决：直线的参数方程、平面方程
以及数量积（\cref{thm:g12:space:planeeq}，
\cref{prop:g12:space:distance}）。全题使用
\cref{exo:g12:space:8}中的立方体：棱长 $1$，$A(0,0,0)$、$B(1,0,0)$、
$D(0,1,0)$、$E(0,0,1)$、$C(1,1,0)$、$F(1,0,1)$、$H(0,1,1)$、
$G(1,1,1)$。

\medskip
\noindent\textbf{第一部分 --- 基本技能。}
\begin{enumerate}
  \item 写出过 $B(1,0,0)$ 且
        法向量为 $(1,1,1)$ 的平面方程，
        过 $A$ 且方向向量为 $(1,1,1)$ 的直线的参数方程，以及
        它们的交点。
  \item 求 $A$ 到该平面的距离
        （\cref{prop:g12:space:distance}）。
  \item 判断向量 $(1,-1,0)$ 与 $(1,1,-2)$ 是否正交？
  \item 判断平面 $x + y + z = 1$ 与
        $2x + 2y + 2z = 5$ 的位置关系，以及过 $A$ 且方向向量为
        $(1,1,0)$ 的直线相对于第一个平面的位置关系？
  \item 写出立方体六条棱
        $[BC], [CD], [DH], [HE], [EF], [FB]$ 的中点，并验证
        这六个点都满足方程
        $x + y + z = \frac32$。
\end{enumerate}

\medskip
\noindent\textbf{第二部分 --- 六边形截面。}
设平面 $P$ 为 $x + y + z = \frac32$。
\begin{enumerate}[resume]
  \item 证明 $P$ 垂直于体对角线
        $(AG)$，且经过立方体的中心。
  \item 计算第 5 题中所列的相邻中点之间的距离，
        沿截面一周：你发现了什么？
  \item 求截面多边形在一个顶点处的角（用
        两条相邻边向量及
        数量积）。由此可得截面是一个
        \emph{正六边形}。
  \item 求该六边形的周长与面积，并将面积与
        立方体一个面的面积作比较。（附带一提：
        单位立方体的最大平面截面是
        $1 \times \sqrt2$ 的对角线矩形，面积为
        $\sqrt2$ --- 正六边形屈居亚军。）
  \item 对角线 $(AG)$ 与 $P$ 的交点在哪里？验证
        该交点是 $[AG]$ 的中点。
  \item 平移截面：描述截面
        $x + y + z = c$ 随 $c$ 由 $0$ 增至 $3$ 的变化 ---
        求 $c = \frac12$ 时的截面（是何种多边形，
        大小如何？），并描述经过
        $c = \frac32$ 时所得六边形的变化过程。
  \item 为什么没有任何平面能把立方体截成一个
        \emph{七}边形？（截面的每条边必须位于
        何处？）
\end{enumerate}

\medskip
\noindent\textbf{第三部分 --- 隐藏的四面体。}
\begin{enumerate}[resume]
  \item 计算角上四面体 $ABDE$ 的体积
        （底面积、高，以及初中分册承认的三分之一公式
        --- 见\cref{exo:g12:space:9}，
        其中已一般地计算过此类体积）。
  \item 四个顶点 $B$、$D$、$E$、$G$：计算全部六个
        两两之间的距离，并断定 $BDEG$ 是一个
        \emph{正}四面体。随后由立方体中减去
        若干角上四面体，求得其体积。
  \item 计算立方体中心到
        平面 $(BDE)$ 的距离。
  \item 金刚石角：计算立方体的
        对角线 $\vect{AG}$ 与 $\vect{BH}$ 的夹角。其
        补角约为 $109.5^\circ$，是金刚石与甲烷中
        化学键之间的夹角 --- 为何碳原子周围的四个电子对
        会选择
        第 14 题中四面体的顶点？
\end{enumerate}

\medskip
\noindent\textbf{第四部分 --- 真正的三维问题。}
\begin{enumerate}[resume]
  \item 求过 $B(1,0,0)$ 且以
        $(0,1,1)$ 为方向向量的直线与六边形所在平面 $P$ 相交于何处。
  \item 证明直线 $(AB)$ 与 $(EG)$ 既不相交
        也不平行：它们是\emph{异面直线}，
        这是平面内不可能出现的现象。
  \item 求点 $B$ 在 $P$ 上的正交投影
        （自 $B$ 沿法向量移动，直至到达平面），并
        验证所得点满足平面 $P$ 的方程。
  \item 总结 --- 空间测绘工具包：
        用于求交点的直线参数方程、用于求角与距离的法向量、
        用于截面的平面方程；以及立方体作为
        试验场，在其中得到了六边形、正
        四面体、金刚石角以及一对异面
        直线。每项各用一句话概括。
\end{enumerate}
\end{problem}
